\section{Drawing a Region of Interest}
\label{sec:region-of-interest}

\subsection{Live hint}

The top of the \emph{Region of Interest} section is a blue-accent
info box that updates live as you change the workflow type. It tells
you exactly which frames need a region:

\begin{itemize}
  \item \emph{Accumulative} $\to$ ``Only frame 1 needs a Region of
    Interest. All later frames are compared against it directly.''
  \item \emph{Incremental + Every Frame} $\to$ ``Frame 1 needs a
    Region of Interest. It is automatically warped forward to each
    later frame.''
  \item \emph{Incremental + Every N Frames} $\to$ ``Draw a Region of
    Interest on frames: 1, 6, 11, 16, \ldots''
  \item \emph{Incremental + Custom Frames} $\to$ Lists your custom
    reference frame indices.
\end{itemize}

\subsection{The Region of Interest toolbar}

Three rows of buttons, from top to bottom:

\subsubsection*{Row 1 --- Shape operations}

\begin{description}[leftmargin=1.2cm,style=nextline]
  \item[+ Add $\blacktriangledown$]
    Opens a popup menu: Polygon / Rectangle / Circle. Pick one, then
    draw on the canvas. The drawn region is added (OR'd) to the
    Region of Interest mask.
  \item[\ding{46} Cut $\blacktriangledown$]
    Same menu, but the drawn region is subtracted (AND-NOT) from the
    Region of Interest mask. Use this to cut holes around features
    you don't want to track (bubbles, voids, damaged areas).
  \item[+ Refine $\blacktriangledown$]
    Paint brush for marking mesh-refinement zones (does \emph{not}
    change the Region of Interest --- it marks areas where the
    quadtree mesh should be subdivided for finer resolution).
\end{description}

\begin{warnbox}
The Refine brush is gated to frame 1 only. On any other frame the
button renders with a dashed border and muted colours; hovering shows
the explanation. This is because the brush paints in frame-1 pixel
coordinates and is automatically warped forward --- painting on a
deformed frame would place refinement zones in the wrong material
points.
\end{warnbox}

\subsubsection*{Row 2 --- Import / Batch Import}

\begin{description}[leftmargin=1.2cm,style=nextline]
  \item[Import]
    Load a grayscale / binary image and use it as the Region of
    Interest for the currently-edited frame. Any pixel above zero
    counts as ``inside''.
  \item[Batch Import]
    Opens a dialog to import several mask files at once and assign
    them to specific frames. Useful when you have pre-segmented
    masks from e.g.\ SAM, thresholding, or a dedicated segmentation
    tool.
\end{description}

\subsubsection*{Row 3 --- Save / Invert / Clear}

\begin{description}[leftmargin=1.2cm,style=nextline]
  \item[Save]
    Writes the current frame's Region of Interest mask to a PNG.
  \item[Invert]
    Flips inside/outside on the current frame.
  \item[Clear]
    Empties the current frame's Region of Interest mask. Does not
    delete other frames' regions.
\end{description}

\subsection{Drawing shapes}

After selecting a shape from Add or Cut:

\begin{description}[leftmargin=1.2cm,style=nextline]
  \item[Polygon]
    Click to place vertices, double-click or press Enter to close.
    Right-click to cancel.
  \item[Rectangle]
    Click-drag to define the bounding box.
  \item[Circle]
    Click the centre, drag outward to set the radius, release to
    commit.
\end{description}

While drawing, the toolbar button for the active mode highlights.
The canvas cursor changes to a crosshair. Click \emph{Esc} or the
already-highlighted button again to cancel.

\subsection{Per-frame editing}

Click the Region column button for any frame to jump to editing that
frame's region. The canvas shows the selected frame's image (in
frame-1 coordinates for accumulative mode, or whichever ref frame
you're editing for incremental). A banner overlays the top of the
canvas:

\begin{center}
  \texttt{EDITING REGION OF INTEREST FOR FRAME 07}
\end{center}

Changes are committed as soon as you finish a shape operation.
